% BAB I Pendahuluan. Sub-bab baku: Pedoman Tabel 52; isi: Subbab 10.3.
\chapter{Pendahuluan}
\label{bab:pendahuluan}

\panduan{BAB I meyakinkan penguji bahwa masalah yang diteliti riil,
penting, dan belum terjawab (Pedoman Subbab 10.3). Struktur BAB I sampai
BAB III identik dengan naskah proposal. Pada naskah akhir, ubah kalimat
rencana menjadi laporan pelaksanaan (Subbab 11.2.1). Penilaian bab ini
masuk ke Sub-CPMK1-CPMK1: masalah empiris mitra harus terurai berbasis
bukti data (Buku Panduan Dosen, Rubrik CPMK1).}

\section{Latar Belakang}
\label{sec:latar-belakang}

\panduan{Susun latar belakang sebagai corong dari konteks luas ke masalah
spesifik (Subbab 10.3.1). Paragraf pertama: konteks bidang secara ringkas.
Paragraf kedua: fenomena lapangan yang Anda observasi di mitra, dengan
bukti dari model proses as-is. Paragraf ketiga: apa yang sudah dikerjakan
peneliti terdahulu. Paragraf penutup: kesenjangan empiris yang dibuktikan
dengan data atau sintesis pustaka pada BAB~\ref{bab:tinjauan}, bukan
dengan kalimat \enquote{belum banyak penelitian}. Hindari kalimat generik
seperti \enquote{di era digital} (Subbab 11.6.1).}

Contoh paragraf konteks dengan satu sitasi: kerangka design science
menempatkan artefak teknologi informasi sebagai objek yang dirancang dan
dievaluasi untuk memecahkan masalah organisasi \cite{hevner2004design}.

Contoh paragraf yang memakai dua sitasi sekaligus: penelitian yang
melibatkan tinjauan pustaka perlu melaporkan prosedur penelusuran dan
seleksi sumber secara terbuka agar dapat diaudit
\cite{hevner2004design,page2021prisma}. Di samping itu, pengukuran kegunaan sistem informasi pada sektor publik secara konsisten mengadopsi instrumen standar yang tervalidasi \cite{refriansyah2025evaluasi,rizky2024evaluasi}.

Contoh paragraf fenomena lapangan: pada kantor desa mitra, tuliskan angka
hasil observasi Anda sendiri, misalnya jumlah permohonan surat per bulan
dan rerata waktu tunggu, lalu rujuk model proses as-is yang Anda susun.
Angka ini harus berasal dari data mitra yang telah diizinkan untuk
dipublikasikan.

\section{Rumusan Masalah}
\label{sec:rumusan-masalah}

\panduan{Tulis satu sampai tiga pertanyaan penelitian operasional yang
dapat dijawab dengan data yang direncanakan pada BAB~\ref{bab:metode}
(Subbab 10.3.2). Rumusan masalah, tujuan, dan metrik keberhasilan tidak
boleh berubah diam-diam antara proposal dan naskah akhir. Setiap
perubahan dicatat dan dijustifikasi pada BAB~\ref{bab:metode}
(Subbab 11.2.1).}

Rumusan masalah penelitian ini adalah sebagai berikut.
\begin{enumerate}
  \item Berapa skor kegunaan aplikasi menurut pengguna sasaran?
  \item Tugas mana yang paling sering gagal diselesaikan responden?
\end{enumerate}

\section{Tujuan Penelitian}
\label{sec:tujuan}

\panduan{Setiap tujuan berpasangan satu lawan satu dengan rumusan masalah
dan memakai kata kerja operasional C4 sampai C6, misalnya membandingkan,
merancang, atau memvalidasi (Subbab 10.3.2, Standar Mutu BAB I butir 4).}

\begin{enumerate}
  \item Mengukur skor kegunaan aplikasi menurut pengguna sasaran.
  \item Menganalisis tugas yang paling sering gagal diselesaikan responden.
\end{enumerate}

\section{Batasan Masalah}
\label{sec:batasan}

\panduan{Batasan menyebut sedikitnya lingkup data, lingkup metode, dan
lingkup waktu (Standar Mutu BAB I butir 6). Batasan yang tegas melindungi
Anda dari ekspektasi yang melampaui kapasitas satu semester.}

\begin{enumerate}
  \item Lingkup data: tuliskan wilayah, periode, dan sumber data.
  \item Lingkup metode: tuliskan metode yang dipakai dan yang dikecualikan.
  \item Lingkup waktu: tuliskan rentang waktu pengambilan data.
\end{enumerate}

\section{Manfaat Penelitian}
\label{sec:manfaat}

\panduan{Pisahkan manfaat teoretis bagi pengembangan keilmuan dan manfaat
praktis bagi mitra atau pengguna. Keduanya harus spesifik
(Standar Mutu BAB I butir 5).}

\begin{enumerate}
  \item Manfaat teoretis: tuliskan kontribusi terhadap kesenjangan yang
        dibuktikan pada BAB~\ref{bab:tinjauan}.
  \item Manfaat praktis: tuliskan keputusan atau perbaikan apa yang dapat
        diambil mitra dari hasil penelitian.
\end{enumerate}

\section{Sistematika Penulisan}
\label{sec:sistematika}

\panduan{Rujuk setiap BAB dengan rujukan silang otomatis, bukan angka
yang diketik manual (Subbab 11.7.1).}

BAB~\ref{bab:pendahuluan} memuat latar belakang, rumusan masalah, tujuan,
batasan, dan manfaat penelitian. BAB~\ref{bab:tinjauan} memuat landasan
teori, sintesis penelitian terdahulu, dan kerangka pemikiran.
BAB~\ref{bab:metode} memuat desain, data, instrumen, tahapan, metrik, dan
etika penelitian. BAB~\ref{bab:hasil} memuat hasil per tujuan penelitian
beserta pembahasannya. BAB~\ref{bab:penutup} memuat kesimpulan dan saran.
